\documentclass[11pt,a4paper]{article}
\usepackage[T1]{fontenc}
\usepackage[utf8]{inputenc}
\usepackage{lmodern}
\usepackage{amsfonts}
\usepackage[margin=28mm]{geometry}
\usepackage{xcolor}
\usepackage{microtype}
\usepackage{parskip}
\definecolor{cimenavy}{HTML}{0D2A42}
\definecolor{cimeblue}{HTML}{164F72}
\definecolor{cimegold}{HTML}{BD8744}
\pagestyle{empty}
\setlength{\parindent}{0pt}
\begin{document}
{\small\bfseries\color{cimegold} CIME SCHOOL 2027}\hfill
{\small\color{cimeblue} Cetraro, Italy | 20-24 September 2027}

\vspace{5mm}
{\color{cimenavy}\rule{\textwidth}{1.2pt}}
\vspace{8mm}

{\fontsize{22}{27}\selectfont\bfseries\color{cimenavy}
Introduction to Weak KAM theory and a generalization of viscosity solutions to metric spaces\par}

\vspace{7mm}
{\large\bfseries Albert Fathi}\par
{\color{cimeblue} ENS de Lyon}

\vspace{10mm}
{\large\bfseries\color{cimegold} Abstract}

Weak KAM theory originally connected the Mather theory of Lagrangian systems and the viscosity solution theory of Hamilton-Jacobi equations. Specifically, it related the Mather theory of a Lagrangian system on a manifold $M$, defined by a Lagrangian function $L:TM \to \mathbb{R}$ (where $TM$ is the tangent bundle of $M$), to the viscosity solutions of the corresponding Hamilton-Jacobi equation. This connection is established when the Lagrangian is convex in the velocity variable.

We will begin by introducing essential concepts from both Lagrangian systems and the viscosity solution theory of Hamilton-Jacobi equations, describing the connection between them.

Given a Lagrangian $L$, adding a suitable constant (called the Ma\~n\'e critical value) allows us to define the Ma\~n\'e potential $\Phi: M \times M \to \mathbb{R}$ as the minimal action required to join two points in $M$ over an arbitrary time interval.

Building on an idea by Antonio Siconolfi, one can demonstrate that key concepts, including Peierls barriers, Aubry sets, viscosity subsolutions and viscosity solutions themselves, can be recovered solely from the Ma\~n\'e potential.

We will show how this fact allows the theory to apply in the more general framework of compact metric spaces, proposing a way to define solutions of the Hamilton-Jacobi equation on general metric spaces.

\vfill
{\color{cimenavy}\rule{\textwidth}{0.5pt}}\\[-1mm]
{\footnotesize\color{cimeblue} Hamilton-Jacobi Equations, Games, and Percolation: Interfaces and Synergies}
\end{document}
